
% BEGIN COMPLETE PARITY AND CONDUCTOR RECEIVERS

\clearpage\begingroup
\part*{Original parity and Pollaczek source bounds}
\newtheorem{paritytheorem}{Theorem}
\newtheorem{parityproposition}[paritytheorem]{Proposition}
\newtheorem{paritylemma}[paritytheorem]{Lemma}
\newcommand{\parityR}{\mathbb R}
\newcommand{\parityC}{\mathbb C}
\newcommand{\parityN}{\mathbb N}
\newcommand{\parityPol}{\mathbb C[x]}
\newcommand{\paritydiag}{\operatorname{diag}}
\newcommand{\parityRan}{\operatorname{Ran}}



\section{The unchanged arithmetic weight and the literature input}

Retain the original quartet, divisor, and completed numerator
\[
 \rho=\tfrac12+\delta+i\gamma,\quad 0<\delta<\tfrac12,\quad\gamma>2,
 \quad h(s)=\prod_{z\in\{\rho,\bar\rho,1-\rho,1-\bar\rho\}}(s-z)^m,
 \quad m\geq1,\quad g=2\xi,
\]
\[
 w_h(y)=\frac{|(g/h)(1/2+iy)|^2}{2\pi},\quad
 m_k=w_h^{*k},\quad k=4l+1\geq9,\quad
 \mu_h=\int_\parityR w_h(y)\,dy.
\]
All multiplicities are the original complete multiplicities. No
existence of an off-line quartet is asserted here. Conjugation
symmetry of $g$ and of the root set of $h$ gives $w_h(-y)=w_h(y)$,
and convolution gives the same symmetry for $m_k$.
The source variable, centre and quotient polynomial remain
\[
 S=c+iy,\quad c=k/2,\quad q=[1+k(m-1)](k+1)^2,
\]
\[
 \chi(S)=\prod_{a,b=0}^k
 [S-c-(2a-k)\delta-i(2b-k)\gamma]^{1+k(m-1)}.
\]
All finite observation formulas below use the original degree
$N\geq q-1$ and actual map $\mathcal A_S=\Lambda\mathcal J_{S,N}$,
with its full physical Taylor unit. Its codomain is its original
observation image with the retained Hermitian structure.

The original proved CAI16--18 envelope is used at its actual scope
\cite{parity:CAI}; the arbitrary-fixed-multiplicity form is HG21 \cite{parity:HG}:
\begin{equation}\label{parity:eq:CAI}
 a_k(1+y^2)^{-M}\rho_\sigma(y)\leq m_k(y)\leq A_k\rho_\sigma(y),
 \quad M=21+4m,
\end{equation}
\[
 \rho_\sigma(y)=\frac{|\Gamma(1/4+iy/2)|^2}{2\pi},\quad
 \int\rho_\sigma=\sqrt{2\pi},\quad\alpha=\pi/2,
\]
\[
 c_\Gamma e^{-\alpha|y|}(1+|y|)^{-1/2}
 \leq\rho_\sigma(y)\leq
 C_\Gamma e^{-\alpha|y|}(1+|y|)^{-1/2},
\]
\[
 c_\Gamma=\sqrt2e^{-7/6},\quad C_\Gamma=\sqrt{2\sqrt5}e^{2/3},
 \quad p=8m-9/2,\quad B_h=42+8m,
\]
\[
 a_k=\frac{c_h\vartheta_h^{k-3}e^{-\alpha(k-3)}(k-2)^{-B_h}}
                   {C_\Gamma2^{B_h/2}},\quad
 A_k=\frac{C_h^{\rm tilt}}{c_\Gamma}k^{p+1}M_\alpha^{k-1},
\]
\[
 \vartheta_h=\int_{-1}^1w_h(y)\,dy,\qquad
 M_\alpha=\int_\parityR e^{\alpha y}w_h(y)\,dy.
\]
The constants $c_h,C_h^{\rm tilt}$ are precisely the established
lower-envelope and compact/tail constants in the complete CAI
provider, not new assumptions. In particular $m_k>0$ and all
moments exist. For complex polynomials $P$ of degree at most $d$,
that provider proves
\begin{equation}\label{parity:eq:CAI-finite}
 \frac{a_k}{\Delta_d}\|P\|_\sigma^2
       \leq\|P\|_{m_k}^2\leq A_k\|P\|_\sigma^2,
 \qquad \Delta_d=[1+4(d+M)^2]^M.
\end{equation}
The Gamma mass has not been divided out.

The original author source used here is A.~I.~Aptekarev,
G.~L\'opez Lagomasino and A.~Mart\'{\i}nez--Finkelshtein,
\emph{Strong asymptotics for the Pollaczek multiple orthogonal
polynomials ensembles},
\href{https://arxiv.org/abs/1410.1261v1}{arXiv:1410.1261v1}
\cite{parity:ALM}.
Its Introduction gives the weights
\begin{equation}\label{parity:eq:literature}
 w_1(x)=\frac1{\sinh(\alpha\sqrt x)},\qquad
 w_2(x)=\frac1{\sqrt x\cosh(\alpha\sqrt x)},\qquad x>0,
\end{equation}
and their positive Stieltjes ratio
\begin{equation}\label{parity:eq:S}
 s(x)=\frac{w_2(x)}{w_1(x)}
     =\frac{\tanh(\alpha\sqrt x)}{\sqrt x}
     =\frac4\pi\sum_{j=0}^{\infty}\frac1{x+(2j+1)^2}.
\end{equation}
The lowercase $s(x)$ in this note denotes this multiplier, not
the original complex source variable $S$. No multiple-orthogonal
asymptotic is assumed for the native weight.

\section{Exact parity-square maps, source images, and masses}

Define the two original transported measures on $(0,\infty)$ by
\begin{equation}\label{parity:eq:nu}
 d\nu_0(x)=\frac{m_k(\sqrt x)}{\sqrt x}\,dx,
 \qquad d\nu_1(x)=\sqrt x\,m_k(\sqrt x)\,dx=x\,d\nu_0(x).
\end{equation}
They have strictly positive densities. Their literal density
ratios against the first literature weight are
\begin{equation}\label{parity:eq:ratios}
 r_0(x)=\frac{m_k(\sqrt x)\sinh(\alpha\sqrt x)}{\sqrt x},
 \qquad r_1(x)=x r_0(x),\qquad d\nu_j=r_jw_1\,dx.
\end{equation}
For every polynomial $f(y)$ there is a unique pair $E,O$ with
$f(y)=E(y^2)+yO(y^2)$. Evenness of $m_k$ makes the cross term in
the squared norm odd. The substitution $x=y^2$ on each half-axis
then proves the exact sesquilinear and norm identities
\begin{equation}\label{parity:eq:parity-norm}
 \|f\|_{m_k}^2=\int_0^\infty|E(x)|^2\,d\nu_0(x)
                    +\int_0^\infty|O(x)|^2\,d\nu_1(x).
\end{equation}
There is no missing factor two: the two half-axes cancel the
$1/2$ in $dy=dx/(2\sqrt x)$.

The same construction is a unitary isomorphism
\[
 L^2(\parityR,m_kdy)\longrightarrow L^2(\nu_0)\oplus L^2(\nu_1),
\]
with
\[
 E(x)=\frac{f(\sqrt x)+f(-\sqrt x)}2,\qquad
 O(x)=\frac{f(\sqrt x)-f(-\sqrt x)}{2\sqrt x}.
\]
These formulas are defined almost everywhere, and their inverse
is $f(y)=E(y^2)+yO(y^2)$ for $y\ne0$. The identity of norms proves
both that the maps are well defined and that they are inverse
isometries.

The exact isometry into the literature's two Hilbert spaces is
\begin{equation}\label{parity:eq:literature-isometry}
 (E,O)\longmapsto
       \left(\sqrt{r_0}\,E,\sqrt{r_1/s}\,O\right)
 \quad\text{in }L^2(w_1dx)\oplus L^2(w_2dx).
\end{equation}
Since all density ratios are positive, division by their square
roots gives its inverse on these full Hilbert spaces. On the
degree-$N$ source its image is exactly
\[
 \sqrt{r_0}\,\parityPol_{\lfloor N/2\rfloor}
       \ \oplus\ \sqrt{r_1/s}\,\parityPol_{\lfloor(N-1)/2\rfloor},
\]
not an unweighted polynomial pair. Here $\parityPol_n$ means degree at
most $n$, and $\parityPol_{-1}=\{0\}$. The common-first-weight version
is $(E,O)\mapsto(\sqrt{r_0}E,\sqrt{r_1}O)$ in two copies of
$L^2(w_1dx)$.

Write $\eta_j=\int y^jw_h(y)\,dy$ and
$\zeta_j=\int y^jm_k(y)\,dy$. Tonelli and the convolution
multinomial formula give the original parity masses
\begin{equation}\label{parity:eq:masses}
 \nu_0((0,\infty))=\zeta_0=\mu_h^k,
 \qquad \nu_1((0,\infty))=\zeta_2=k\eta_2\mu_h^{k-1}.
\end{equation}
The cross terms in the second formula vanish because $w_h$ is
even. The literature masses are different fixed numbers:
\begin{equation}\label{parity:eq:raw-moments}
 \tau_j:=\int_0^\infty x^j w_1(x)\,dx
 =\frac{4(2j+1)!}{\alpha^{2j+2}}
        (1-2^{-2j-2})\zeta(2j+2),\qquad \tau_0=2,
\end{equation}
\[
 \int_0^\infty w_2(x)\,dx=2.
\]
For the first formula use $x=y^2$ and the positive expansion
$1/\sinh(\alpha y)=2\sum_{j\geq0}e^{-(2j+1)\alpha y}$,
then integrate each power of $y$. Its constant case uses
$\sum_{j\geq0}(2j+1)^{-2}=\pi^2/8$, also obtained below from
\eqref{parity:eq:S} at zero. For the second mass, the same substitution
gives $2\int_0^\infty\operatorname{sech}(\alpha y)dy=\pi/\alpha=2$;
an antiderivative of $\operatorname{sech} u$ is
$\arctan(\sinh u)$. None of these masses is set equal to one.

For the exact original source $F(S)=\sum_{j=0}^Na_jS^j$, define
\[
 T_{rj}=\binom jr c^{j-r}i^r\quad(r\leq j),\qquad T_{rj}=0\quad(r>j).
\]
Let $\Pi_N$ reorder the coefficients in powers of $y$ into even
powers first and odd powers second, without rescaling any entry.
The parity coefficient map is the invertible matrix
\begin{equation}\label{parity:eq:source-map}
 L_N=\Pi_NT,\qquad (e,o)=L_Na,
\end{equation}
where explicitly
\[
 e_r=(-1)^r\sum_{j=2r}^N\binom j{2r}c^{j-2r}a_j,
 \quad
 o_r=i(-1)^r\sum_{j=2r+1}^N\binom j{2r+1}c^{j-2r-1}a_j.
\]
Thus the full original degree-$N$ source maps onto precisely the
displayed pair of polynomial spaces. An original constrained
source subspace $W$ maps to $L_NW$, not to a substituted full pair.

Put $n_0=\lfloor N/2\rfloor$, $n_1=\lfloor(N-1)/2\rfloor$ and
\begin{equation}\label{parity:eq:parity-grams}
 H_{j,n}=\left[\int_0^\infty x^{a+b}\,d\nu_j(x)\right]_{a,b=0}^n,
 \qquad H=\paritydiag(H_{0,n_0},H_{1,n_1}).
\end{equation}
The matrices are positive definite because nonzero polynomials
cannot vanish on a set of full measure for a positive density.
They retain the literal moments
$(H_{j,n})_{ab}=\zeta_{2(a+b+j)}$. Integration gives
\begin{equation}\label{parity:eq:original-gram}
 H_{S,N}^{\rm ar}=L_N^*HL_N.
\end{equation}
The actual observation matrix in these coordinates is
\begin{equation}\label{parity:eq:observation}
A=\mathcal A_SL_N^{-1}.
\end{equation}
The exact original observation metric and minimizing lift are
\begin{equation}\label{parity:eq:Q}
 Q(H)=(AH^{-1}A^*)^{-1},\qquad
 a_{\min}=L_N^{-1}H^{-1}A^*Q(H)u.
\end{equation}
Indeed $A$ is onto its stated observation image, so $AH^{-1}A^*$
is strictly positive. The vector $z_0=H^{-1}A^*Q(H)u$ satisfies
$Az_0=u$, and for $Az=0$ the cross term is
$z^*Hz_0=(Az)^*Q(H)u=0$. Every other lift has energy
$u^*Q(H)u+z^*Hz$, proving the formula and uniqueness. The
congruence \eqref{parity:eq:original-gram} returns the result to the
unchanged source coefficients.

There is also an exact transport of the original quotient, not
only of its source norm. Put
\[
 u_a=(2a-k)\delta,\qquad v_b=(2b-k)\gamma,\qquad
 e_k=1+k(m-1),
\]
and define the monic polynomial
\begin{equation}\label{parity:eq:Xi}
 \Xi(x)=\prod_{a=0}^k\prod_{b=(k+1)/2}^{k}
                 [x-(v_b-iu_a)^2]^{e_k}.
\end{equation}
Its degree is $q/2$. The opposite grid point of $(a,b)$ is
$(k-a,k-b)$, and exactly one point of each opposite pair has
$v_b>0$. That pair's factors after $S=c+iy$ are
\[
 (iy-u_a-iv_b)(iy+u_a+iv_b)
          =-[y^2-(v_b-iu_a)^2].
\]
Since $k$ is odd, $q=e_k(k+1)^2$ is divisible by four. The number
$q/2$ of paired factors with multiplicity is therefore even,
so the full sign is $(-1)^{q/2}=1$. Consequently
\begin{equation}\label{parity:eq:chi-square}
 \chi(c+iy)=\Xi(y^2)
\end{equation}
with exactly the original leading coefficient. Conjugating the
roots in \eqref{parity:eq:Xi} is the permutation $a\mapsto k-a$, so
$\Xi$ also has real coefficients.

Let $\mathcal E_x=\parityC[x]/(\Xi)$. Substitution and parity give the
explicit algebra isomorphism followed by vector-space coordinates
\begin{equation}\label{parity:eq:quotient-parity}
 \parityC[S]/(\chi)\longrightarrow\parityC[y]/(\Xi(y^2))
                  \longrightarrow\mathcal E_x\oplus\mathcal E_x,
 \qquad [F]\longmapsto([E],[O]).
\end{equation}
The second arrow has multiplication
$(E,O)(E',O')=(EE'+xOO',EO'+OE')$ modulo $\Xi$ in each component.
To verify its kernel, divide $E$ and $O$ separately by $\Xi$.
The remaining even and odd polynomials have degrees less than
$q$ in $y$, so their sum is divisible by the monic degree-$q$
polynomial $\Xi(y^2)$ only when both remainders are zero. This
also proves surjectivity and identifies the inverse through the
unique remainder $E(y^2)+yO(y^2)$ and $y=(S-c)/i$.
Multiplication by the original variable $S$ is precisely
\begin{equation}\label{parity:eq:S-block}
 \begin{pmatrix}cI&iM_x\\ iI&cI\end{pmatrix}
 \quad\text{on }\mathcal E_x\oplus\mathcal E_x.
\end{equation}
An original full Taylor unit $\mathfrak u\in\parityC[S]/(\chi)$,
transported as $(U_0,U_1)$, acts by the unchanged block matrix
\begin{equation}\label{parity:eq:unit-block}
 \begin{pmatrix}M_{U_0}&M_{xU_1}\\M_{U_1}&M_{U_0}\end{pmatrix}.
\end{equation}
Its inverse is the transported multiplication by
$\mathfrak u^{-1}$; the unit has not been removed. Thus the
original class and observation maps can equivalently be
transported through \eqref{parity:eq:quotient-parity} and
\eqref{parity:eq:unit-block}. Formula \eqref{parity:eq:observation} is their
same matrix composition before choosing these quotient
coordinates. At $N\geq q-1$, both parity source degree bounds
are at least $q/2-1$, so both residue-class components are
attained by the original degree-$N$ source.

\section{Both native orthogonal families and all Christoffel constants}

Let $P_d(y)$ be the real monic orthogonal polynomial of degree $d$
for $m_k(y)dy$ and let $\omega_d=\int P_d^2m_k>0$. Existence and
uniqueness follow by solving the positive definite moment system.
Evenness and uniqueness give $P_d(-y)=(-1)^dP_d(y)$, so define
\[
 P_{2n}(y)=E_n(y^2),\quad P_{2n+1}(y)=yO_n(y^2),
 \quad \kappa_n=\omega_{2n},\quad\lambda_n=\omega_{2n+1}.
\]
Equation \eqref{parity:eq:parity-norm} shows that $E_n$ and $O_n$ are
respectively the real monic orthogonal families for $\nu_0$ and
$\nu_1$, with norms $\kappa_n$ and $\lambda_n$.

The monic orthogonal polynomial in the original complex source
coordinate is not obtained by dropping the phase. It is exactly
\begin{equation}\label{parity:eq:S-polynomial-phase}
 \mathsf P_d(S)=i^dP_d((S-c)/i),\qquad
 \mathsf P_d(c+iy)=i^dP_d(y),\qquad
 \|\mathsf P_d\|_{m_k}^2=\omega_d.
\end{equation}
Its leading coefficient is $i^di^{-d}=1$; the lower-degree
orthogonality follows by the same invertible substitution used
in \eqref{parity:eq:source-map}. Thus uniqueness proves this exact
dictionary. Its parity pair is $((-1)^nE_n,0)$ at degree $2n$
and $(0,i(-1)^nO_n)$ at degree $2n+1$. These phases remain in
the original boundary classes and observation entries.

Orthogonality gives the full original recurrence
\[
 yP_d=P_{d+1}+\beta_dP_{d-1},\qquad
 \beta_d=\omega_d/\omega_{d-1}>0\quad(d\geq1),\quad \beta_0=0.
\]
To prove it, expand $yP_d-P_{d+1}$ in the lower orthogonal basis.
Its coefficients below $d-1$ vanish by transferring $y$ to the
lower polynomial. The coefficient of $P_d$ vanishes by parity,
and the coefficient of $P_{d-1}$ is
$\langle P_d,yP_{d-1}\rangle/\omega_{d-1}=\omega_d/\omega_{d-1}$.
In the original $S$ coordinate the same identity reads
$(S-c)\mathsf P_d=\mathsf P_{d+1}-\beta_d\mathsf P_{d-1}$,
since $i^{d+1}/i^{d-1}=-1$; its sign is therefore retained too.

Set
\begin{equation}\label{parity:eq:ad}
 a_n=\frac{\lambda_n}{\kappa_n}=\beta_{2n+1}>0,
 \qquad d_n=\frac{\kappa_n}{\lambda_{n-1}}=\beta_{2n}>0\ (n\geq1),
 \qquad d_0=0.
\end{equation}
Separating the even and odd instances of that recurrence gives
the exact polynomial identities
\begin{equation}\label{parity:eq:Christoffel}
 xO_n=E_{n+1}+a_nE_n,\qquad
 E_n=O_n+d_nO_{n-1},\quad O_{-1}=0.
\end{equation}
Evaluating the first at zero proves, starting from $E_0=1$,
\begin{equation}\label{parity:eq:E0}
 E_n(0)=(-1)^n\prod_{j=0}^{n-1}a_j\ne0,
 \qquad \frac{E_{n+1}(0)}{E_n(0)}=-a_n.
\end{equation}
Therefore the exact Christoffel formula is
\begin{equation}\label{parity:eq:Christoffel-explicit}
 O_n(x)=\frac{E_{n+1}(x)-[E_{n+1}(0)/E_n(0)]E_n(x)}x
        =\frac{E_{n+1}(x)+a_nE_n(x)}x,
 \quad \lambda_n=a_n\kappa_n.
\end{equation}
The numerator has a zero at zero, so this is a polynomial identity,
including the endpoint. The complete odd zero value is
\begin{equation}\label{parity:eq:O0}
 O_n(0)=(-1)^n\sum_{j=0}^n
     \left(\prod_{r=j+1}^nd_r\right)
     \left(\prod_{r=0}^{j-1}a_r\right)
 =E_{n+1}'(0)+a_nE_n'(0).
\end{equation}
The sum follows by iterating $O_n(0)=E_n(0)-d_nO_{n-1}(0)$;
the derivative formula follows from \eqref{parity:eq:Christoffel-explicit}.
Every empty product equals one, so this covers $n=0$.

Substitution of the two identities in \eqref{parity:eq:Christoffel} into
each other yields both complete monic recurrences
\begin{align}
 xE_n&=E_{n+1}+(a_n+d_n)E_n+a_{n-1}d_nE_{n-1},\label{parity:eq:E-rec}\\
 xO_n&=O_{n+1}+(a_n+d_{n+1})O_n+a_nd_nO_{n-1}.\label{parity:eq:O-rec}
\end{align}
The last term of the first is zero at $n=0$ and needs no $a_{-1}$.
The norm ratios are exactly
$\kappa_n/\kappa_{n-1}=a_{n-1}d_n$ and
$\lambda_n/\lambda_{n-1}=a_nd_n$.
For example the orthonormal version of the first cross relation is
\[
 x\frac{O_n}{\sqrt{\lambda_n}}
  =\sqrt{d_{n+1}}\frac{E_{n+1}}{\sqrt{\kappa_{n+1}}}
       +\sqrt{a_n}\frac{E_n}{\sqrt{\kappa_n}},
\]
and the second is
\[
 \frac{E_n}{\sqrt{\kappa_n}}
 =\sqrt{a_n}\frac{O_n}{\sqrt{\lambda_n}}
       +\sqrt{d_n}\frac{O_{n-1}}{\sqrt{\lambda_{n-1}}}.
\]
These equalities are between functions in the displayed different
weighted spaces; they explicitly give the matrix of the identity
map or $x$-multiplication on every finite polynomial domain.

\section{The exact multiple-orthogonality receiver and its residuals}

The raw paper has one polynomial $P$ subject to moments against
both $w_1$ and $w_2$. The native common-multiplier Nikishin pair is
\[
 (d\nu_0,\ s\,d\nu_0)=r_0(x)(w_1(x),w_2(x))\,dx;
\]
the native odd measure is instead $d\nu_1=x\,d\nu_0$.
The literal maps between their functionals, for every polynomial
$P$ and nonnegative integer $j$, are
\begin{align}
 \int x^jP\,d\nu_0&=\int x^j(r_0P)w_1\,dx,\label{parity:eq:functional0}\\
 \int x^jP\,s\,d\nu_0&=\int x^j(r_0P)w_2\,dx,\label{parity:eq:functionalS}\\
 \int x^jP\,d\nu_1&=\int x^j\left(\frac{r_1}{s}P\right)w_2\,dx
                    =\int x^{j+1}P\,d\nu_0.\label{parity:eq:functional1}
\end{align}
Thus each finite raw moment operator receives the explicit
multiplied source space in these formulas, not the raw polynomial
space unless that multiplier is retained. These maps are the
functional counterparts of the square-root isometry
\eqref{parity:eq:literature-isometry}.

Here is a nonzero residual computation within the native common
multiplier pair. The $n$ zeros of $E_n$ are simple and lie in
$(0,\infty)$. To prove this, multiply $E_n$ by the product of its
distinct sign-changing zeros in that interval. If fewer than $n$
such zeros existed, this product would have degree less than $n$
and its product with $E_n$ would have constant nonzero sign except
at finitely many points of the positive interval. Positivity of
$\nu_0$ would contradict orthogonality. Thus there are $n$ such
zeros, which exhaust the degree and establish the assertion.
In particular $E_n(-a)$ has sign $(-1)^n$ for every $a>0$.

For $a>0$, divide $E_n(x)-E_n(-a)$ by $x+a$. The quotient has
degree $n-1$, so orthogonality gives the exact positive-resolvent
identity
\begin{equation}\label{parity:eq:residual-atom}
 \int\frac{E_n(x)}{x+a}\,d\nu_0(x)
  =\frac1{E_n(-a)}\int\frac{E_n(x)^2}{x+a}\,d\nu_0(x).
\end{equation}
For $0\leq j<n$, division of $x^j$ by $x+a$ similarly gives
\begin{equation}\label{parity:eq:residual}
 \int x^jE_n(x)s(x)\,d\nu_0(x)
 =\frac4\pi\sum_{r=0}^{\infty}
 \frac{[-(2r+1)^2]^j}{E_n(-(2r+1)^2)}
 \int\frac{E_n(x)^2}{x+(2r+1)^2}\,d\nu_0(x).
\end{equation}
To justify the sum before division, dominate by
$\int|x^jE_n(x)|s(x)\,d\nu_0<\infty$; the atomwise integrals are
absolutely summable since $s$ is bounded. The divisions preserve
each integral exactly. Every term on the resulting right side
has the same strict sign $(-1)^{n+j}$, so that is the sign of the
nonzero residual. In particular $E_n$ does not satisfy even the
zeroth second-weight condition of the native common-multiplier
Nikishin system. This is an evaluated residual, not just a
declaration that two orthogonalities differ. For $n=0$ the
zeroth residual is directly positive.

There is also a completely finite lower bound. Cauchy--Schwarz
with the factors $(x+1)^{1/2}E_n$ and $(x+1)^{-1/2}E_n$ yields
\[
 \int\frac{E_n^2}{x+1}\,d\nu_0
 \geq\frac{\kappa_n^2}{\int(x+1)E_n^2\,d\nu_0}
 =\frac{\kappa_n}{1+a_n+d_n},
\]
where \eqref{parity:eq:E-rec} computes the final moment. The first
positive atom in \eqref{parity:eq:residual} therefore gives
\begin{equation}\label{parity:eq:residual-lower}
 \left|\int E_n s\,d\nu_0\right|
 \geq\frac{4\kappa_n}{\pi(1+a_n+d_n)|E_n(-1)|}>0.
\end{equation}
All its constants belong to the original native parity family.

\section{Two positive resolvent series, with all constants proved}

We give a direct proof of \eqref{parity:eq:S} and the reciprocal identity
that will supply finite native bounds. On the interval
$[0,\alpha]$, $\alpha=\pi/2$, the operator $-d^2/dt^2+x$ with
Dirichlet condition at zero and Neumann condition at $\alpha$
has eigenvalues $x+(2j+1)^2$ and normalized eigenfunctions
$\sqrt{2/\alpha}\sin((2j+1)t)$. Their completeness is the ordinary
sine Fourier expansion obtained by odd reflection at zero and
even reflection at $\alpha$. The Green kernel is
\[
 G_x(u,v)=\frac{\sinh(\sqrt x\min(u,v))
                  \cosh(\sqrt x(\alpha-\max(u,v)))}
                    {\sqrt x\cosh(\alpha\sqrt x)}.
\]
Its formula follows by solving the homogeneous differential
equation on either side of $u=v$, imposing the two boundary
conditions, continuity, and the unit derivative jump. Its
eigenfunction expansion is uniformly absolutely convergent on
the closed square since its terms are bounded by
$(2/\alpha)/(x+(2j+1)^2)$. The spectral resolvent and the displayed
Green kernel agree in $L^2$ and hence by continuity everywhere.
At $u=v=\alpha$ this gives \eqref{parity:eq:S}; at $x=0$, continuity
gives $s(0)=\alpha$ and the odd reciprocal-square sum used above.

With Neumann conditions at both endpoints, the constant mode
contributes $1/(\alpha x)$ and the positive modes have
eigenvalues $x+4j^2$, $j\geq1$. The normalized endpoint values
have squares $2/\alpha$. The analogous Green kernel has endpoint
value $\coth(\alpha\sqrt x)/\sqrt x$. Multiplication by $x$ gives
the exact reciprocal series
\begin{equation}\label{parity:eq:reciprocal}
 \boxed{\frac1{s(x)}=\sqrt x\coth(\alpha\sqrt x)
  =\frac2\pi+\frac4\pi\sum_{j=1}^{\infty}\frac{x}{x+4j^2}.}
\end{equation}
The same uniformly convergent kernel argument justifies the
expansion for $x>0$. Both sides extend at zero to $2/\pi$.
Every summand is nonnegative. These identities are used on the
original axis $x>0$; no degree-dependent rescaling is made.

\section{Finite Stieltjes compression and native Gram approximation}

The following formulas apply to each original parity block with
$\nu=\nu_j$ and degree $n=n_j$. Put $v_n(x)=(1,x,\ldots,x^n)^T$,
\[
 H=\int v_nv_n^*\,d\nu,\quad
 X_H=\int x v_nv_n^*\,d\nu,\quad
 K=\int v_nv_n^*s\,d\nu,\quad
 K_x=\int x v_nv_n^*s\,d\nu.
\]
All four are finite positive definite matrices, and for the
zeroth parity block $X_H=H_{1,n}$ exactly. For $a>0$ define
$R(a)=\int v_nv_n^*/(x+a)\,d\nu$. The block Gram matrix of
$(x+a)^{1/2}v_n$ and $(x+a)^{-1/2}v_n$ is positive. Its Schur
complement gives
\begin{equation}\label{parity:eq:resolvent-compression}
 R(a)\succeq H(X_H+aH)^{-1}H.
\end{equation}
The positive series \eqref{parity:eq:S} therefore proves
\begin{equation}\label{parity:eq:finite-S}
 H^{1/2}s(H^{-1/2}X_HH^{-1/2})H^{1/2}
       \preceq K\preceq\alpha H.
\end{equation}
The matrix function in the lower bound is defined by the positive
spectral decomposition of the finite matrix. The summation is
justified because each inverse is at most $a^{-1}I$ and the
reciprocal-square series converges. This is a finite-moment
bound using the exact Christoffel moment matrix, not a replacement
of the full multiplication operator by its compression.

More directly useful for the unchanged original Gram is
\eqref{parity:eq:reciprocal}. Define the actual second-weight tilt
$d\widetilde\nu=s\,d\nu=r_jw_2\,dx$ and, for an integer $J\geq1$,
\begin{equation}\label{parity:eq:H-trunc}
 H^{[J]}=\frac2\pi K+\frac4\pi\sum_{j=1}^J
       \int\frac{x}{x+4j^2}v_nv_n^*\,d\widetilde\nu.
\end{equation}
Every summand is positive and $H^{[J]}$ is strictly positive
already from its first term. Integrating \eqref{parity:eq:reciprocal}
and using monotone convergence in every quadratic form gives
$H^{[J]}\uparrow H$. Since
$x/(x+4j^2)\leq x/(4j^2)$ and
$\sum_{j>J}j^{-2}\leq\int_J^\infty t^{-2}dt=1/J$, the full
matrix error bound is
\begin{equation}\label{parity:eq:H-error}
 \boxed{0\preceq H-H^{[J]}\preceq\frac1{\pi J}K_x
                                      \preceq\frac1{2J}X_H.}
\end{equation}
This is a certified approximation of the original native Gram,
not just a representation of an auxiliary tilted metric.

Only moments and negative-axis resolvent integrals are required
for a finite truncation. If
$\widetilde\mu_d=\int x^d\,d\widetilde\nu$ and
$\widetilde F(a)=\int(x+a)^{-1}\,d\widetilde\nu$, polynomial
division gives, for $d\geq1$,
\begin{equation}\label{parity:eq:resolvent-moments}
 \int\frac{x^d}{x+a}\,d\widetilde\nu
 =\sum_{r=0}^{d-1}(-a)^r\widetilde\mu_{d-1-r}
                  +(-a)^d\widetilde F(a).
\end{equation}
For $d=0$ the integral is $\widetilde F(a)$. This proves the
finite evaluation map, including every polynomial-division sign.
The positive integral form in \eqref{parity:eq:H-trunc} remains available
when cancellation in \eqref{parity:eq:resolvent-moments} would be
numerically inconvenient.

Apply these constructions to both blocks of
\eqref{parity:eq:parity-grams}. Denote their direct sums by $H^{[J]}$,
$K_x$ and $X_H$. The original source error is exactly the
congruence
\[
 0\preceq H_{S,N}^{\rm ar}-L_N^*H^{[J]}L_N
       \preceq\frac1{\pi J}L_N^*K_xL_N.
\]
For the unchanged onto observation $A$ in \eqref{parity:eq:observation},
the minimum-energy proof of \eqref{parity:eq:Q} shows that the map
$G\mapsto Q(G)=(AG^{-1}A^*)^{-1}$ is order-preserving on positive
source matrices: taking the minimum over the same set $Az=u$
preserves each quadratic-form inequality. Therefore
\begin{equation}\label{parity:eq:Q-bracket}
 \boxed{Q(H^{[J]})\preceq Q(H)\preceq
 Q\left(H^{[J]}+\frac1{\pi J}K_x\right).}
\end{equation}
Both endpoints converge to the original $Q(H)$: the finite
matrices converge to $H$, their positive inverses converge by
$B^{-1}-C^{-1}=B^{-1}(C-B)C^{-1}$, and the same holds after
the onto observation congruence. This proves a finite certified
source and observation calculation with the original classes and
Taylor unit untouched.

There is an entirely explicit degree-only bound for the error.
The original Gamma orthonormal recurrence has off-diagonal
$b_j=\sqrt{(j+1)(j+1/2)}\leq j+1$, so its two weighted shifts
give $\|yP\|_\sigma\leq2(N+1)\|P\|_\sigma$ on degree at most $N$.
Combining this with \eqref{parity:eq:CAI-finite} yields
\begin{equation}\label{parity:eq:L-bound}
 \|yP\|_{m_k}^2\leq\mathcal L_N\|P\|_{m_k}^2,
 \qquad \mathcal L_N=
       \frac{4A_k\Delta_N(N+1)^2}{a_k}.
\end{equation}
Under the parity map this is exactly $X_H\preceq\mathcal L_NH$.
Consequently for $J>\mathcal L_N/2$, with
$\varepsilon_J=\mathcal L_N/(2J)<1$, one has
\begin{equation}\label{parity:eq:relative}
 (1-\varepsilon_J)H\preceq H^{[J]}\preceq H,
 \qquad
 Q(H^{[J]})\preceq Q(H)
       \preceq(1-\varepsilon_J)^{-1}Q(H^{[J]}).
\end{equation}
This bound is generally coarse but is finite, positive and
explicit in the original CAI constants and degree. It does not
assume the desired arithmetic estimate.

\section{A second fully explicit bound using raw Pollaczek moments}

The literal ratio $r_0$ also admits a useful finite comparison
with the raw first literature weight. Put $L=M+1$ and
\begin{equation}\label{parity:eq:ratio-constants}
 c_0=\frac{a_kc_\Gamma\alpha}{2^{1/4}(1+2\alpha)},
 \qquad C_0=A_kC_\Gamma\alpha.
\end{equation}
Then for every $x>0$,
\begin{equation}\label{parity:eq:ratio-envelope}
 c_0(1+x)^{-L}\leq r_0(x)\leq C_0.
\end{equation}
To prove this, write $y=\sqrt x$ and retain the exact factor
\[
 e^{-\alpha y}\frac{\sinh(\alpha y)}y
    =\frac{1-e^{-2\alpha y}}{2y}.
\]
It is at most $\alpha$ and at least
$\alpha/(1+2\alpha y)$, since $1-e^{-u}\geq u/(1+u)$ for $u\geq0$.
The latter inequality follows from $e^u\geq1+u$.
Substituting the two CAI envelopes gives the upper bound and the
lower expression
\[
 \frac{a_kc_\Gamma\alpha(1+y^2)^{-M}}
                   {(1+2\alpha y)(1+y)^{1/2}}.
\]
Use $1+2\alpha y\leq(1+2\alpha)\sqrt{1+y^2}$ and
$(1+y)^{1/2}\leq2^{1/4}(1+y^2)^{1/4}$, followed by
$(1+y^2)^{-M-3/4}\geq(1+y^2)^{-M-1}$, to prove
\eqref{parity:eq:ratio-envelope} with the stated constants.

For parity $j=0,1$ use the raw measure
$d\tau^{(j)}=x^jw_1(x)dx$ and the degree $n=n_j$. Define the
explicit moment matrices
\[
 U_{j,n}=[\tau_{a+b+j}]_{a,b=0}^n,\qquad
 V_{j,n}=\left[\sum_{r=0}^L\binom Lr\tau_{a+b+j+r}\right]_{a,b=0}^n.
\]
All entries are the explicit values \eqref{parity:eq:raw-moments}.
The measures satisfy $d\nu_j=r_0\,d\tau^{(j)}$, so the same
block-Gram Schur-complement argument used in
\eqref{parity:eq:resolvent-compression}, now with weight $(1+x)^L$,
proves the actual native bound
\begin{equation}\label{parity:eq:raw-matrix-bound}
 \boxed{c_0 U_{j,n}V_{j,n}^{-1}U_{j,n}
             \preceq H_{j,n}\preceq C_0U_{j,n}.}
\end{equation}
Explicitly, the block matrix with blocks $V_{j,n},U_{j,n}$ and
$\int v_nv_n^*(1+x)^{-L}d\tau^{(j)}$ is a Gram matrix, giving
the lower rational-weight bound; \eqref{parity:eq:ratio-envelope}
then gives the displayed native inequality.
Taking block direct sums, followed by the unchanged congruence
$L_N$ or the unchanged observation $A$, transports these bounds
to the original source or to $Q(H)$, respectively. In particular,
if $G_-$ and $G_+$ are these two block bounds, then
\[
 Q(G_-)\preceq Q(H)\preceq Q(G_+).
\]
This gives a fully explicit finite comparison from raw literature
moments to the actual observation metric, with every density
ratio paid for. No raw multiple-orthogonal polynomial is being
used as a native orthogonal polynomial.

The same ratio envelope proves a typed map in the reverse
direction for the paper's actual multiple orthogonal polynomial
$P_{(n_1,n_2)}$. The function
\[
 f(x)=P_{(n_1,n_2)}(x)/r_0(x)
\]
belongs to $L^2(\nu_0)$, since its squared norm is
$\int |P|^2r_0^{-1}w_1dx$ and the inverse ratio is bounded by
$c_0^{-1}(1+x)^L$. Equations
\eqref{parity:eq:functional0}--\eqref{parity:eq:functionalS} give exactly
\[
 \int x^jf\,d\nu_0=0\ (0\leq j<n_1),\qquad
 \int x^jf\,s\,d\nu_0=0\ (0\leq j<n_2).
\]
This is an explicit receiver of the raw multiple-orthogonality
relations in the original even Hilbert sector. Its source image
is $r_0^{-1}\parityPol$, not the original polynomial source. For the
paper's diagonal rescaling
$Q_n(x)=(4n^2)^{-2n}P_{(n,n)}(4n^2x)$, the corresponding original
even-axis function is exactly
\[
 y\longmapsto\frac{(4n^2)^{2n}Q_n(y^2/(4n^2))}{r_0(y^2)}.
\]
This formula retains the complete coordinate change and multiplier;
no asymptotic statement about native $E_n$ or $O_n$ follows from
the paper's asymptotic theorem by omitting them.

\section{What the finite estimates do and do not complete}

The parity isometries, Christoffel maps and Gram congruences are
exact representations of the original objects. The nonzero
residual \eqref{parity:eq:residual-lower}, the rational truncation and
error bound \eqref{parity:eq:H-error}, the original observation brackets
\eqref{parity:eq:Q-bracket}--\eqref{parity:eq:relative}, and the explicit
raw-moment comparison \eqref{parity:eq:raw-matrix-bound} are additional
finite calculations and inequalities. Their constants and
receiving maps have all been evaluated above.

They preserve the original complex observed cross entries
\[
 B_i=(\Lambda a_i)^*Q(H)\Lambda v,
 \qquad
 \Phi_{B,N}=\frac{2\Re(\overline{B_1}B_2)}{\omega_{N,k}^{\rm ar}},
\]
with the actual boundary classes, terminal class, and physical
Taylor unit. The bounds certify source/observation approximation
and do not supply an evaluated sign, phase, or required
asymptotic of this last signed product. No such conclusion, and
no imported multiple-orthogonal asymptotic, is claimed.

\begin{thebibliography}{9}
\bibitem{parity:ALM} A. I. Aptekarev, G. L\'opez Lagomasino and
A. Mart\'{\i}nez--Finkelshtein,
\emph{Strong asymptotics for the Pollaczek multiple orthogonal
polynomials ensembles}, arXiv:1410.1261v1 (2014), Introduction:
the two unscaled weights and their positive Stieltjes ratio.
\url{https://arxiv.org/abs/1410.1261v1}.
The original author LaTeX is the source used; its multiple-orthogonal
asymptotic theorem is not assumed for the arithmetic weight.
\bibitem{parity:CAI} Split-Zero programme, \emph{The complete original
action relative to its arithmetic total}, CAI16--18 and its complete original
convolution providers. This is the programme calculation supplying
the actual arithmetic density comparison, not a human-literature
attribution for the Pollaczek weights.
\url{https://github.com/KokunoYumeto/zeta-function-research-reader/blob/2412f4aaa07c3fb19052cc0e537e93c2db9b0b1d/workbenches/splitzero-tandem/continuations/20260920-heat-literature-observed-class/providers/CAI_COMPLETE_PREREQUISITE.tex}.
\bibitem{parity:HG} Split-Zero programme, \emph{Heat growth and original
currents}, HG21 and its displayed proof, retaining every fixed
primary multiplicity and the original Gamma mass.
\url{https://github.com/KokunoYumeto/zeta-function-research-reader/blob/2412f4aaa07c3fb19052cc0e537e93c2db9b0b1d/workbenches/splitzero-tandem/continuations/20260920-heat-literature-observed-class/providers/HEAT_GROWTH_AND_ORIGINAL_CURRENTS.tex}.
The parity recurrences, Green-kernel expansions, finite Schur
inequalities and receiving maps are proved explicitly in this note;
no omitted asymptotic argument is attributed to these programme
providers.
\end{thebibliography}


\endgroup

\clearpage\begingroup
\part*{Growing-degree truncation and the full conductor}
\newcommand{\nativecutC}{\mathbb C}
\newcommand{\nativecutR}{\mathbb R}

This note carries the exact parity/Pollaczek receiver into the original
arithmetic multiplication estimate and the original conductor quotient.
The human input is the unscaled weight ratio of Aptekarev, L\'opez
Lagomasino and Mart\'{\i}nez--Finkelshtein \cite{nativecut:ALM}; all transport to
the arithmetic measure is proved in the complete companion \cite{nativecut:PP}.
The stronger truncation length below uses the already proved original
CAI21--25 estimate \cite{nativecut:CAI}, not a new comparison measure. No native
resolvent integral or signed phase is assigned an uncomputed value.

\section{Original objects and the proved finite truncation}
Retain the complete quartet divisor $h$, $g=2\xi$, the even source
$w_h(y)=|(g/h)(1/2+iy)|^2/(2\pi)$, and $m_k=w_h^{*k}$ with
$k=4\ell+1\ge9$, as in the companion. Their masses remain
$\mu_h^k$. For $S=c+iy$, $c=k/2$, the literal polynomial coordinate
map is $L_N=\Pi_NT_N$, where
$(T_N)_{rj}=\binom jr c^{j-r}i^r$ for $r\le j$ and $\Pi_N$ orders
the even powers first. The original source splits as
$f(y)=E(y^2)+yO(y^2)$ with the exact measures
\[
 d\nu_0(x)=m_k(\sqrt x)x^{-1/2}dx,\qquad d\nu_1=x\,d\nu_0.
\]
Let $H$ be their direct-sum monomial Gram at degrees
$\lfloor N/2\rfloor,\lfloor(N-1)/2\rfloor$. Then the original
$S$-coordinate Gram is $H_S=L_N^*HL_N$.

The literature multiplier and its exact reciprocal are
\[
 s(x)=\frac{\tanh(\pi\sqrt x/2)}{\sqrt x},\qquad
 \frac1{s(x)}=\frac2\pi+\frac4\pi\sum_{j\ge1}\frac{x}{x+4j^2}.
\]
Both identities, including convergence and constants, are proved from
the one-dimensional Green kernels in \cite{nativecut:PP}. For each parity block,
put $v(x)=(1,x,\ldots,x^n)^T$ and retain the actual tilt
$d\widetilde\nu=s\,d\nu$. Define
\[
 H^{[J]}=\frac2\pi\int vv^*d\widetilde\nu
 +\frac4\pi\sum_{j=1}^J\int\frac{x}{x+4j^2}vv^*d\widetilde\nu,
 \qquad X_H=\int xvv^*d\nu,
\]
and use direct sums for the two blocks. Positivity of every omitted
term and $\sum_{j>J}j^{-2}\le1/J$ give
\begin{equation}\tag{TR1}\label{nativecut:eq:tail}
 0\preceq H-H^{[J]}\preceq\frac1{\pi J}\int xvv^*s\,d\nu
                                  \preceq\frac1{2J}X_H.
\end{equation}
The last step uses $s(x)\le\pi/2$. All three integrals are actual
native integrals, not the raw Pollaczek moments.

\section{The original tensor bound sharpens the truncation length}
Here are the complete constants needed from CAI21--25 \cite{nativecut:CAI}:
\[
 \alpha=\pi/2,\ M=21+4m,\ \nu_*=\log(5/3),\ d_* =\alpha-\nu_*>0,
 \quad C_\sigma=\frac4{\sqrt{15}\sqrt{2\pi}},
\]
\[
 \begin{gathered}
 H_h=\frac{64C_\sigma(1+4M^2)^M}{d_*^2}
             \max_{b=3,4,5}\frac{M_\alpha^b}{a_b},\\
 T_h=\frac2{d_*}\{\log16+2M+3+\max(0,\log H_h)\},\qquad
 C_{\rm mult}(h)=\sqrt{T_h^2+1}.
 \end{gathered}
\]
The constants $a_b,M_\alpha$ are the original displayed constants
of the companion's CAI envelope, with the convolution order replaced
by $b$. CAI21 proves
\[
 \int_{|y|>R}y^2|P(y)|^2m_b(y)dy
 \le H_h(n+1)^{2M+2}16^ne^{-d_*R/2}\|P\|_{m_b}^2
 \quad(\deg P\le n,\ b=3,4,5).
\]
At $R=T_h(n+1)$ its coefficient is at most one, using
$\log(n+1)\le n$; inside the interval $y^2\le R^2$. Addition gives
$\|yP\|_{m_b}\le C_{\rm mult}(h)(n+1)\|P\|_{m_b}$.
Thus the tail result is used only at its established block orders.

Write $r_k=\lfloor k/3\rfloor$ and partition $k$ into $r_k$
summands $b_i\in\{3,4,5\}$, replacing one three by four or five
when necessary. Convolution proves the exact isometry
\[
 f(y)\longmapsto f(y_1+\cdots+y_{r_k})
 \quad\hbox{into }\bigotimes_iL^2(m_{b_i}(y_i)dy_i).
\]
Its image is the sum-dependent subspace and its mass is still
$\mu_h^k$. In each original orthonormal polynomial basis,
multiplication has zero diagonal by evenness and off-diagonal
coefficients $0<a_n^{(i)}\le C_{\rm mult}(h)n$: take the degree-$n$
coefficient in the preceding multiplication bound at degree $n-1$.

Compress the sum of multiplication operators to the product basis
of total degree at most $N+1$. At a vertex of degree at most $N$
its absolute row sum is at most $C_{\rm mult}(h)(2N+r_k)$; on the
outer face it has only downward entries and row sum at most
$C_{\rm mult}(h)(N+1)\le C_{\rm mult}(h)(2N+r_k)$. The symmetric
matrix norm is at most its maximal absolute row sum, by pairing
off-diagonal terms and using $2|uv|\le |u|^2+|v|^2$. Multiplication
of a degree-$N$ sum-dependent polynomial has degree at most $N+1$,
so no output is lost. Returning through the isometry proves
\begin{equation}\tag{TR2}\label{nativecut:eq:native}
 \|yP\|_{m_k}^2\le C_{\rm mult}(h)^2(2N+r_k)^2\|P\|_{m_k}^2.
\end{equation}
The companion also gives the full Gamma comparison
$\|yP\|_{m_k}^2\le[4A_k\Delta_N(N+1)^2/a_k]\|P\|_{m_k}^2$,
where $\Delta_N=[1+4(N+M)^2]^M$. Therefore retain both bounds as
\begin{equation}\tag{TR3}\label{nativecut:eq:eps}
 L_N^\sharp=\min\left\{\frac{4A_k\Delta_N(N+1)^2}{a_k},
                     C_{\rm mult}(h)^2(2N+r_k)^2\right\},
 \qquad \epsilon_{N,J}=\frac{L_N^\sharp}{2J}.
\end{equation}
Under the exact parity map, TR2 is $X_H\preceq L_N^\sharp H$.
For the actual integers $J>L_N^\sharp/2$, TR1 yields
\begin{equation}\tag{TR4}\label{nativecut:eq:relative}
 (1-\epsilon_{N,J})H\preceq H^{[J]}\preceq H.
\end{equation}
The sharper bound is quadratic in $N+k$ for fixed $h$. It does not
require an asymptotic estimate for the native orthogonal polynomials.

\section{Actual quotient metrics and signed endpoint errors}
Let $A$ be the original onto observation expressed in the parity
coordinates: $A=\Lambda J_{S,N}L_N^{-1}$. For a positive source
matrix $G$, the original attained minimum is
$Q_A(G)=(AG^{-1}A^*)^{-1}$. Its minimizing lift is
$G^{-1}A^*Q_A(G)u$; every other lift differs by a vector in $\ker A$
orthogonal to it for $G$. This proves the formula, and minimising
TR4 over the same fibre gives
\begin{equation}\tag{TR5}\label{nativecut:eq:Q}
 Q_{N,J}:=Q_A(H^{[J]})\preceq Q_N:=Q_A(H)
           \preceq(1-\epsilon_{N,J})^{-1}Q_{N,J}.
\end{equation}
Restriction to any specified original polynomial source subspace
preserves TR4 and the identical argument applies to its specified
onto observation. No different admitted relation space is substituted.

Write $d_N=\dim\operatorname{im}A$. The eigenvalues of
$Q_{N,J}^{-1/2}Q_NQ_{N,J}^{-1/2}$ lie between one and
$(1-\epsilon_{N,J})^{-1}$. Their product gives
\begin{equation}\tag{TR6}\label{nativecut:eq:det}
 0\le\log\frac{\det Q_N}{\det Q_{N,J}}
                 \le-d_N\log(1-\epsilon_{N,J}).
\end{equation}
This calculation proves an inequality for determinants in their
original frames; the square-root conjugation does not replace those
frames or norms. The empty determinant is one and satisfies TR6.

For the original endpoints $N_1,N_2,N_3,N_4$, retain
\[
 F=\log\frac{\det Q_{N_1}\det Q_{N_2}}
                  {\det Q_{N_3}\det Q_{N_4}},\qquad
 F_{\mathbf J}=\log\frac{\det Q_{N_1,J_1}\det Q_{N_2,J_2}}
                  {\det Q_{N_3,J_3}\det Q_{N_4,J_4}}.
\]
Put $e_i=-d_{N_i}\log(1-\epsilon_{N_i,J_i})$. Each individual
logarithmic error is nonnegative, so the complete signed enclosure is
\begin{equation}\tag{TR7}\label{nativecut:eq:four}
 -e_3-e_4\le F-F_{\mathbf J}\le e_1+e_2.
\end{equation}
For every specified tolerance $t>0$, choose actual integers
\begin{equation}\tag{TR8}\label{nativecut:eq:J}
 J_i\ge\left\lceil L_{N_i}^\sharp\max\{1,2d_{N_i}/t\}\right\rceil.
\end{equation}
Then $\epsilon_{N_i,J_i}\le1/2$ and
$-\log(1-\epsilon)\le2\epsilon$ give $e_i\le t/2$, hence
$|F-F_{\mathbf J}|\le t$. When $d_{N_i}=0$ its contribution is
zero and the same integer choice remains valid. At the packet
windows $N_i=O(q)$, $d_{N_i}\le q$ and fixed $h$, this uses
$O_h(q^3/t)$ resolvent terms for $0<t\le1$. Taking $t=1/k$ gives
a vanishing original four-endpoint truncation error. This is a
truncation-length bound, not a claim about numerical conditioning,
total complexity, or certification cost for the native integrals.
It does not evaluate the signed return or its limiting value.

\section{The same certified bound reaches the full conductor}
On the simple-quartet conductor domain, EQ1--21 \cite{nativecut:EQ} proves
$C=\bar C\Lambda$ and the full attained conductor metric
\[
 T_N=(\bar C Q_N^{-1}\bar C^*)^{-1},\qquad
 R_N=Q_N^{-1}\bar C^*T_N,
\]
with the same original source and with $\bar C$ onto. Apply the
same fixed-fibre minimum to TR5. Defining
$T_{N,J}=(\bar C Q_{N,J}^{-1}\bar C^*)^{-1}$ gives
\begin{equation}\tag{TR9}\label{nativecut:eq:conductor}
 T_{N,J}\preceq T_N\preceq
                    (1-\epsilon_{N,J})^{-1}T_{N,J}.
\end{equation}
The quotient identity also follows by direct substitution:
$T_{N,J}=(\bar CA(H^{[J]})^{-1}A^*\bar C^*)^{-1}$.
Thus the same truncation is applied to the original conductor
map, rather than to a separately chosen four-plane metric.
Restricting by its original cubic inclusion $I_W$ and coordinate
map $\Psi$ gives the identical relative bracket for
$G_N=\Psi^*I_W^*T_NI_W\Psi$. EQ16--19's exact pair-flip matrix
$B_J=O(I-D_J)O^{-1}E$ is independent of this metric. For every
original label vector $\alpha$ its actual defect energy obeys
\begin{equation}\tag{TR10}\label{nativecut:eq:energy}
 \alpha^*B_J^*G_{N,j}B_J\alpha
 \le\alpha^*B_J^*G_NB_J\alpha
 \le(1-\epsilon_{N,j})^{-1}\alpha^*B_J^*G_{N,j}B_J\alpha.
\end{equation}
Here the lower-case $j$ is the truncation length, distinct from
the two-element sign-flip set $J$; $G_{N,j}$ uses $T_{N,j}$.
The rank-two mixing map, all state labels, the Taylor unit and
the complete conductor kernel are unchanged. TR10 certifies its
positive defect energy, not the sign of a different projected
complex product. The terminal-class residual in EQ20--21 is
also retained; none of these bounds sets it to zero.

\begin{thebibliography}{9}
\bibitem{nativecut:ALM} A. I. Aptekarev, G. L\'opez Lagomasino and
A. Mart\'{\i}nez--Finkelshtein, \emph{Strong asymptotics for the
Pollaczek multiple orthogonal polynomials ensembles},
arXiv:1410.1261v1 (2014), Introduction, unscaled weights and
Stieltjes ratio. \url{https://arxiv.org/abs/1410.1261v1}.
The original author LaTeX was read in full by the integrating task.
\bibitem{nativecut:PP} Split-Zero programme, \emph{Exact native parity transport
into the Pollaczek pair and certified finite source--observation bounds},
complete companion \texttt{PARITY\_POLLACZEK\_NATIVE\_RECEIVER.tex},
equations \texttt{eq:reciprocal}, \texttt{eq:H-error} and
\texttt{eq:Q-bracket}, with their full proofs. This is the actual
arithmetic receiving calculation, not the author paper's asymptotics.
\bibitem{nativecut:CAI} Split-Zero programme, \emph{The complete original action
relative to its arithmetic total}, CAI21--25, in the complete provider
\texttt{CAI\_COMPLETE\_PREREQUISITE.tex}.
\url{https://github.com/KokunoYumeto/zeta-function-research-reader/blob/2412f4aaa07c3fb19052cc0e537e93c2db9b0b1d/workbenches/splitzero-tandem/continuations/20260920-heat-literature-observed-class/providers/CAI_COMPLETE_PREREQUISITE.tex}.
The original complete proof supplies the tail constant used above;
the total-degree transport and every new endpoint bound are proved here.
\bibitem{nativecut:EQ} Split-Zero programme, \emph{The eight signed states
through the original conductor minimum}, complete companion
\texttt{EIGHT\_STATE\_ORIGINAL\_MINIMUM.tex}, EQ1--21; independent
full derivation \texttt{EIGHT\_STATE\_QUOTIENT\_REDERIVATION.tex},
IQ1--29. Human provenance of the upstream factor-map construction
and exact source hashes are retained in those sources and the manifest.
\end{thebibliography}

\endgroup

\clearpage\begingroup
\part*{Eight signed states in the original observation minimum}
\newcommand{\eightminC}{\mathbb C}
\newcommand{\eightminPp}{\mathcal P}
\newcommand{\eightminim}{\operatorname{im}}
\newcommand{\eightminrank}{\operatorname{rank}}
\newcommand{\eightmindiag}{\operatorname{diag}}



This calculation receives the nonlinear ES--Fable construction through
the original finite conductor and the actual arithmetic observation
quotient. It evaluates the inherited eight-state minimum, its exact
correction to the four-dimensional source minimum, and the failure of
evaluation alone to retain the signed monodromy action. Every minimum
below is evaluated as a finite matrix of the unchanged original density.
No asymptotic allocation or arithmetic projected sign is assigned by
these finite identities. The human source of the underlying factor-map
construction is identified in Tao's exposition \cite{eightmin:Tao}; the marked
four-dimensional polynomial and its conductor receiver are the local
programme results \cite{eightmin:FC}, rather than results attributed to Tao.

\section{Original spaces and the actual factor through observation}
Use the simple-packet domain of OB1--3 \cite{eightmin:OB}:
\[
 \begin{gathered}
 \rho=\tfrac12+\delta+i\gamma,\quad0<\delta<\tfrac12,\quad\gamma>2,\\
 k\equiv1\pmod4,\ k\ge9,\quad j=k-8,\quad
 q=(k+1)^2,\quad q_-=(k-7)^2.
 \end{gathered}
\]
The fixed original period is in FC17's admitted five-orbit domain:
the five original transformed quadrics are pairwise nonassociate.
FC17 proves on this domain that the original symbol is nonzero and
its first nonzero moment has order at most eighty. This is the same
domain used by the receiving conductor; it is not an assertion about
an exceptional period outside that domain. All coefficients remain
unchanged. Write
\[
 \chi_k(S)=\prod_{a,b=0}^k
 [S-k/2-(2a-k)\delta-i(2b-k)\gamma],\qquad
 E_k=\eightminC[S]/(\chi_k),
\]
and define $\chi_j,E_j$ by the same expression with $k$ replaced by $j$.
The actual source density is $m_k(y)=w_h^{*k}(y)$ at $S=k/2+iy$,
where $h$ and $w_h$ retain OB1/LR2's complete original factors.
For every $N\ge q-1$, let $H_N$ be its positive Gram in the ordered
monomials $1,S,\ldots,S^N$. Positivity follows since a nonzero
polynomial has positive squared integral against this positive density.

Write $J_N:\eightminPp_N\to E_k$ for reduction modulo $\chi_k$ and
$\Lambda_k:E_k\to B_k$ for the original observation onto its image.
Set $\mathsf A_N=\Lambda_kJ_N$. OB3 proves the exact map
\[
 C:E_k\longrightarrow E_j,\quad C[P]=[T_AP],\qquad
 \ker\Lambda_k\subset\ker C.                                      \tag{EQ1}
\]
Thus the formula $\bar C(\Lambda_kx)=Cx$ is well-defined: two lifts
differ in $\ker\Lambda_k$, so have equal $C$-images. It is linear
and unique because $\Lambda_k$ is onto. In particular, the actual
polynomial map $\mathsf C_N=CJ_N$ satisfies
\[
 \bar C\mathsf A_N=\mathsf C_N.                                    \tag{EQ2}
\]
This constructs the needed receiving map. It does not assume that
$\Lambda_k$ descends through the different quotient
$\eightminPp_{v+3}/\eightminPp_{v-1}$ used in the four-plane construction.

Retain $v=\operatorname{ord}_0E_A\le80$ and the original moments
$\mu_n$. The finite-shift conductor has
\[
 T_A S^n=\sum_{r=0}^n\binom nr\mu_{n-r}(S')^r,
 \quad\mu_0=\cdots=\mu_{v-1}=0,\quad\mu_v\ne0.                    \tag{EQ3}
\]
The triangular map $\eightminPp_N\to\eightminPp_{N-v}$ is onto: its nonzero
diagonal from $S^{v+r}$ to $(S')^r$ is
$\binom{v+r}{r}\mu_v$. Moreover
$q-q_-=16k-48\ge96>80$, so $N-v\ge q_--1$.
Reduction onto $E_j$ is therefore onto. Consequently $\mathsf C_N$
and $\bar C$ are onto, with all of their original kernels retained.

The observation metric, the complete conductor quotient metric, and
their attained lifts are
\[
 \begin{split}
 Q_N&=(\mathsf A_NH_N^{-1}\mathsf A_N^*)^{-1},\\
 T_N&=(\mathsf C_NH_N^{-1}\mathsf C_N^*)^{-1}
       =(\bar C Q_N^{-1}\bar C^*)^{-1},\\
 L_N&=H_N^{-1}\mathsf C_N^*T_N:E_j\longrightarrow\eightminPp_N,\\
 R_N&=Q_N^{-1}\bar C^*T_N:E_j\longrightarrow B_k.
 \end{split}                                                       \tag{EQ4}
\]
All inverses exist by the surjectivity just proved. Direct substitution
gives $\mathsf C_NL_N=I$, $\bar CR_N=I$, and
$\mathsf A_NL_N=R_N$. If $x=L_Nu+d$ and $\mathsf C_Nd=0$, then
$d^*H_NL_Nu=d^*\mathsf C_N^*T_Nu=0$. Hence
\[
 \|x\|_{H_N}^2=u^*T_Nu+d^*H_Nd.                                  \tag{EQ5}
\]
This proves the full original minimum and uniqueness of its lift.
The identical calculation in $B_k$ proves
$R_N^*Q_NR_N=T_N$. Thus EQ4 is equality of successive actual
minima, not an asserted comparison of differently chosen norms.

The available Jacobi reconstruction \cite{eightmin:JC} evaluates every entry
used here from the original one-factor arithmetic transform. With its
unchanged matrices $D_N,C_N^{\rm Jac}$ one has
$H_N=D_N^*C_N^{\rm Jac}D_N$. Substitution in EQ4 keeps the physical
Taylor unit in $\mathsf A_N$ and retains the full kernel. No Gamma
replacement of $m_k$ is made in this substitution.

\section{Exact correction from the four-plane to the complete minimum}
Let $W=\eightminPp_3$ in the unchanged target variable $S'$ and
$I_W:W\to E_j$ be reduction modulo $\chi_j$. Since $q_-\ge4$,
this map is injective. Put
\[
 T_{W,N}=I_W^*T_NI_W.                                               \tag{EQ6}
\]
These are the norms of the complete original conductor lifts of
the actual cubic target values. The four-plane source minimum can
also be calculated exactly in the same arithmetic density.

Partition the principal Gram on $\eightminPp_{v+3}$ into the degrees below
$v$ and degrees $v,\ldots,v+3$:
\[
 H_{v+3}=\begin{pmatrix}H_{LL}&H_{LR}\\H_{RL}&H_{RR}\end{pmatrix},
 \quad S_v=H_{RR}-H_{RL}H_{LL}^{-1}H_{LR},
 \quad (B_v)_{rt}=\begin{cases}\binom{v+t}{r}\mu_{v+t-r},&r\le t,\\
 0,&r>t,\end{cases}                                                \tag{EQ7}
\]
for $0\le r,t\le3$. For $v=0$, the empty lower block is omitted
and $S_v=H_{RR}$. The diagonal of $B_v$ is
$\binom{v+t}{v}\mu_v$, so $B_v$ is invertible without altering
any moment. The unique minimum in $\eightminPp_{v+3}$ mapping to a cubic
with coefficient column $w$ is
\[
 L_{\rm low}w=\operatorname{incl}_{v+3,N}
 \begin{pmatrix}-H_{LL}^{-1}H_{LR}\\I_4\end{pmatrix}B_v^{-1}w,
 \qquad G_{\rm low}=B_v^{-*}S_vB_v^{-1}.                            \tag{EQ8}
\]
Indeed EQ3 forces the four upper coefficients to be $B_v^{-1}w$;
completing the square in the remaining lower coefficients proves
EQ8 and leaves all lower cross terms present. Define
\[
 D_N^{\rm lift}=L_{\rm low}-L_NI_W.
\]
Both lifts have the same $\mathsf C_N$-image, so
$\mathsf C_ND_N^{\rm lift}=0$. Applying EQ5 column by column gives
the exact positive correction
\[
 \boxed{G_{\rm low}-T_{W,N}
       =(D_N^{\rm lift})^*H_ND_N^{\rm lift}\succeq0.}                \tag{EQ9}
\]
In particular equality holds precisely when the two lift matrices
are identical. This answers the metric-identification question with
an explicit matrix of original moments and original maps. The four
source dimensions have not been used to remove the higher-degree
relations from the complete minimum.

\section{All eight states and their inherited observation minimum}
Use the actual four roots and their two factor signs in SE1--9
\cite{eightmin:FC}. Write an eight-state vector uniquely as
$\iota_+\alpha+\iota_-\beta$, retaining the factors $1/2$ in
$\iota_\pm e_r=(f_{r,+}\pm f_{r,-})/2$. The evaluated vector is
$z=E\alpha+O\beta$ in the original ES coordinates. SE7--8 proves
$O$ invertible and $\eightminrank E=3$. Retain the exact original maps
$Y:\eightminC^4\to W$, $\Psi=YK^{-1}:\eightminC^4\to W$ from FC32.
Thus the original target evaluation is $\Psi z$; no coordinate
column is silently declared orthonormal.

The actual lift into the original observation space is now defined,
rather than inferred from the existence of an abstract observation:
\[
 \mathcal O_N(\alpha,\beta)=R_NI_W\Psi(E\alpha+O\beta)\in B_k.
                                                                    \tag{EQ10}
\]
It satisfies $\bar C\mathcal O_N=I_W\Psi(E\alpha+O\beta)$.
Both $R_NI_W\Psi$ and $O$ are injective, so
\[
 \ker\mathcal O_N=\{(\alpha,-O^{-1}E\alpha):\alpha\in\eightminC^4\}.
                                                                    \tag{EQ11}
\]
This identifies exactly the four-dimensional graph kernel in the
original observation space, including the part invisible under sign
averaging.

Retain both labels and evaluation with the map
$\mathcal J(\alpha,\beta)=(Y\alpha,\Psi(E\alpha+O\beta))$.
The full original conductor minimum induces on its two target copies
the orthogonal sum $T_{W,N}\oplus T_{W,N}$. This is a specified
pullback from two copies of the actual minimum, not a previously
given probability measure on the eight abstract states. Define
\[
 H=Y^*T_{W,N}Y,\qquad G=\Psi^*T_{W,N}\Psi,
 \qquad
 \mathbb H_N=\begin{pmatrix}
 H+E^*GE&E^*GO\\O^*GE&O^*GO
 \end{pmatrix}.                                                     \tag{EQ12}
\]
The same proof applies to the original fixed Gamma target form
of FC56 by replacing $T_{W,N}$ with that actual form; the two
applications do not assert equality between those forms.
The map $(\alpha,\beta)\mapsto(\alpha,z)$ is invertible, with
$\beta=O^{-1}(z-E\alpha)$. Hence its full squared norm is exactly
\[
 \alpha^*H\alpha+z^*Gz.                                             \tag{EQ13}
\]
For the observation value $R_NI_W\Psi z_0$, its fibre is exactly
$z=z_0$. Minimisation in EQ13 therefore gives
\[
 \boxed{\min_{\mathcal O_N f=R_NI_W\Psi z_0}\|f\|_{\mathbb H_N}^2
       =z_0^*Gz_0=\|R_NI_W\Psi z_0\|_{Q_N}^2,\quad
       f_{\min}=\iota_-O^{-1}z_0.}                                  \tag{EQ14}
\]
This evaluates the inherited quotient with no additional norm loss
from the four graph-kernel directions. Its relation to the restricted
source cost is exactly EQ9, not an assumption that that correction
vanishes. Indeed $G$ equals
$\Psi^*G_{\rm low}\Psi-\Psi^*(D_N^{\rm lift})^*H_ND_N^{\rm lift}\Psi$.

For completeness the entire inverse of EQ12 is
\[
 \mathbb H_N^{-1}=
 \begin{pmatrix}
 H^{-1}&-H^{-1}E^*O^{-*}\\
 -O^{-1}EH^{-1}&O^{-1}(G^{-1}+EH^{-1}E^*)O^{-*}
 \end{pmatrix}.                                                     \tag{EQ15}
\]
Multiplication by EQ12 verifies this formula; alternatively it
follows by inverting the two triangular maps in EQ13.
In particular $[E\ O]\mathbb H_N^{-1}[E\ O]^*=G^{-1}$.
For any further actual linear observation $A z$ onto its image,
the inherited minimum is $(AG^{-1}A^*)^{-1}$, with unique minimizing
state $(\alpha,\beta)=(0,O^{-1}G^{-1}A^*(AG^{-1}A^*)^{-1}b)$.
This last identity is a fully evaluated finite extension; it does not
postulate an unidentified arithmetic observation map.

\section{The nonzero action defect on those exact fibres}
SM1--14 \cite{eightmin:FC} proves that the actual monodromy contains every
even change of the four factor signs. For a two-element subset
$J\subset\{1,2,3,4\}$ write $D_J$ for its diagonal sign change.
In the coordinates $(\alpha,z)$ of EQ13 its action is precisely
\[
 (\alpha,z)\longmapsto(\alpha,B_J\alpha+L_Jz),\quad
 L_J=OD_JO^{-1},\qquad B_J=O(I-D_J)O^{-1}E.                           \tag{EQ16}
\]
This is the restriction of the complete action SE16, derived directly
by substituting $z=E\alpha+O\beta$ and $\beta\mapsto D_J\beta$.
The canonical section of evaluation is $s(z)=(0,O^{-1}z)$.
On that section $g_Js(z)=s(L_Jz)$. On a general state in the
same evaluation fibre, the extra term is exactly $B_J\alpha$.
The action is therefore not determined by the original observed
value alone whenever this term survives the observation.

Here its size is also an original metric expression. Put
$f_\alpha=(\alpha,-O^{-1}E\alpha)\in\ker\mathcal O_N$. Then
\[
 \mathcal O_N(g_J f_\alpha)=R_NI_W\Psi B_J\alpha,\quad
 \|f_\alpha\|_{\mathbb H_N}^2=\alpha^*H\alpha,\quad
 \|\mathcal O_N(g_J f_\alpha)\|_{Q_N}^2
       =\alpha^*B_J^*GB_J\alpha.                                   \tag{EQ17}
\]
Thus its exact squared norm on the forgotten kernel is
$\lambda_{\max}(H^{-1/2}B_J^*GB_JH^{-1/2})$. There is no replacement
of this expression by the determinant of the nonlinear map.

Every pair flip has a rank-two defect in EQ17. Indeed SE1 gives
the first row of $O$ as $(a_1,a_2,a_3,a_4)$ with all $a_r\ne0$;
SE8 gives $\eightminim E=\{x:x_1=0\}$. Therefore
\[
 \eightminim(O^{-1}E)=\{x:\textstyle\sum_r a_rx_r=0\}.                     \tag{EQ18}
\]
Projection of this hyperplane onto either prescribed pair of
coordinates is onto: choose an omitted coordinate, whose coefficient
is nonzero, to solve the one equation. Consequently
\[
 \eightminim B_J=O\operatorname{span}\{e_r:r\in J\},\quad\eightminrank B_J=2,
 \qquad\sum_{|J|=2}\eightminim B_J=\eightminC^4.                                  \tag{EQ19}
\]
All equalities use the literal original $a_r$, without fixing their
signs. Since $R_NI_W\Psi$ is injective, each of the six pair flips
has rank-two failure to preserve $\ker\mathcal O_N$.
Thus none induces an action on its evaluation quotient.

More generally, for every nonzero linear observation $A$ on the
four evaluated coordinates, at least one actual pair flip has
$AB_J\ne0$. Otherwise EQ19 would imply $A=0$. Choose $\alpha$
with $AB_J\alpha\ne0$. The states $0$ and $f_\alpha$ have the same
original observation zero, but their transformed observations differ
by $AB_J\alpha$. This proves non-descent even for set maps, not only
linear candidates. The full exact receiver is EQ16 on both retained
coordinates, and its missing term after forgetting the labels is
the explicitly evaluated EQ17. Restriction to the chosen section
is a different, stated operation: $L_{J_1}L_{J_2}=OD_{J_1}D_{J_2}O^{-1}$
does give a representation of the sign subgroup there, but it does
not repair the representative dependence just exhibited.

\section{Scope of the receiving result}
The nonlinear example is therefore carried through an actual map
into $B_k$, with its entire inherited quotient metric calculated.
The finite observation values supplied here form precisely the
four-dimensional image $R_NI_W W$. No additional arithmetic
eigenclass has been inserted into it. In particular the lower
terminal cardinal $v_-$ in OB3 has degree $q_--1$: for $k=9$ this
is three, while for $k\ge13$ it is at least 35. Reduction modulo
$\chi_j$ has a unique representative of degree below $q_-$, so
in the latter range $v_-$ does not belong to $I_WW$. The exact
bridge to every original target class nevertheless remains EQ4;
for the terminal class it gives the attained lift
$R_N(a_{88}v_-)$ of the literal conductor image
$Cv_+=a_{88}v_-$. This statement retains the original scalar and
does not assert equality with the original class $\Lambda_kv_+$.
Their difference lies in $\ker\bar C$, and its exact norm identity is
\[
 \|\Lambda_kv_+\|_{Q_N}^2
 =|a_{88}|^2v_-^*T_Nv_-
   +\|\Lambda_kv_+-R_N(a_{88}v_-)\|_{Q_N}^2.                        \tag{EQ20}
\]
The cross term vanishes by the same orthogonality as EQ5, now in
$B_k$. This also gives the actual attained original polynomial.
Define $M_N=H_N^{-1}\mathsf A_N^*Q_N$, so
$\mathsf A_NM_N=I$, $M_N^*H_NM_N=Q_N$, and EQ2--4 gives
$M_NR_N=L_N$. Therefore
\[
 M_N\Lambda_kv_+
 =L_N(a_{88}v_-)+M_N(\Lambda_kv_+-R_N(a_{88}v_-)),\qquad
 \mathsf C_NM_N(\Lambda_kv_+-R_N(a_{88}v_-))=0.                       \tag{EQ21}
\]
Its two summands are $H_N$-orthogonal by EQ5 and have precisely
the two energies in EQ20. Thus the result supplies both the exact
full-class return and the complete residual, without mistaking the
eight-state four-plane for the entire growing packet.

\begin{thebibliography}{9}
\bibitem{eightmin:Tao} Terence Tao, \emph{A digestion of the Jacobian conjecture
counterexample}, 21 July 2026.
\url{https://terrytao.wordpress.com/2026/07/21/a-digestion-of-the-jacobian-conjecture-counterexample/}.
Human exposition of the Alp\"oge--Fable construction; it does not
contain the programme's marked four-dimensional receiver or EQ1--20.
\bibitem{eightmin:FC} The Clankers, ES--Fable marked polynomial construction,
canonical ES reader version 69, theorem
\texttt{explicit-four-time-Jacobian-collision}; programme continuation,
\emph{Fable to original conductor}, FC26--71, GF1--28, SE1--18,
SM1--19. The full exact source versions and reading coverage are in
the accompanying source receipt. Published predecessor:
\url{https://github.com/KokunoYumeto/zeta-function-research-reader/tree/369f6ea9e0312bdf348ca9efd294735efc99c51d}.
\bibitem{eightmin:OB} Programme original observed-class return, OB1--3;
original WCF finite-shift conductor and AC coefficient providers.
\url{https://github.com/KokunoYumeto/zeta-function-research-reader/tree/2412f4aaa07c3fb19052cc0e537e93c2db9b0b1d/workbenches/splitzero-tandem/continuations/20260920-heat-literature-observed-class}.
\bibitem{eightmin:JC} Programme, \emph{The actual arithmetic multiplier in
Gamma Jacobi coordinates}, 20 September 2026, equations
\texttt{eq:sourcegram}, \texttt{eq:quotient}; applying
Franti\v{s}ek \v{S}tampach and Pavel \v{S}\v{t}ov\'{i}\v{c}ek,
\emph{Spectral representation of some weighted Hankel matrices and
orthogonal polynomials from the Askey scheme}, arXiv:1811.05820v1.
The finite minimum proofs EQ4--20 are given here in full.
\end{thebibliography}

\endgroup

\clearpage\begingroup
\part*{Independent full minimum and monodromy-defect proof}
\newcommand{\indminC}{\mathbb C}
\newcommand{\indminim}{\operatorname{im}}
\newcommand{\indmindiag}{\operatorname{diag}}
\newcommand{\indminrank}{\operatorname{rank}}
\newcommand{\indminSpan}{\operatorname{span}}
\newcommand{\indminid}{\operatorname{id}}



\section{Exact scope and the retained source maps}
This note derives finite-dimensional identities from the actual eight-state
maps SE1--18 in \texttt{FABLE\_TO\_ORIGINAL\_CONDUCTOR.tex}, read together
with the explicit monodromy generators SM1--19 of that source.  All quotient
minima below retain the supplied positive Hermitian forms.  No inner product
on the eight abstract labels is assigned independently of the stated original
two-copy map.  In particular, the factors $1/2$ in the following embeddings
are retained:
\[
 \iota_+e_j=(f_{j,+}+f_{j,-})/2,\qquad
 \iota_-e_j=(f_{j,+}-f_{j,-})/2.                       \tag{IQ1}
\]
The actual matrices $E=(e_1\ \cdots\ e_4)$ and $O=(o_1\ \cdots\ o_4)$
have $\indminrank E=3$, $\det O\ne0$, and
\[
 F=O^{-1}E,\qquad
 \indminim F=\{x\in\indminC^4:a_1x_1+\cdots+a_4x_4=0\},\quad a_j\ne0. \tag{IQ2}
\]
These are exactly SE7, SE8 and SE18, with the original branches $a_j$.
The four-plane target maps $Y,\Psi:\indminC^4\longrightarrow W$ are invertible.
For $h=\iota_+\alpha+\iota_-\beta$ write
$\mathcal E h=E\alpha+O\beta$ and set
\[
 z=E\alpha+O\beta,\qquad
 \mathcal J_W h=(Y\alpha,\Psi z),\qquad
 H=Y^*\langle\ ,\ \rangle_WY,\quad
 G=\Psi^*\langle\ ,\ \rangle_W\Psi.                  \tag{IQ3}
\]
Here the expressions for $H,G$ denote the Gram matrices in the displayed
coordinate bases.  Thus $H,G$ are positive definite, with every original
cross entry and mass present.  The exact inverse and inherited norm are
\[
 h=\iota_+\alpha+\iota_-O^{-1}(z-E\alpha),\qquad
 \|h\|_{\mathcal J_W}^2=\alpha^*H\alpha+z^*Gz.         \tag{IQ4}
\]

Let $\Lambda:W\to Z$ be a specified linear observation, with $Z=\indminim\Lambda$;
choose any fixed coordinate basis of this image.  Put $A=\Lambda\Psi$.
Then $A$ is an $m\times4$ matrix of rank $m=\dim Z$, and the eight-state
observation is
\[
 \mathcal R h=\Lambda\Psi(E\alpha+O\beta)=Az.          \tag{IQ5}
\]
This is a finite observation calculation, not an assertion that a different
programme observation descends through a polynomial quotient.  In particular,
an observation defined on a larger original polynomial space cannot be applied
to $\mathcal P_{v+3}/\mathcal P_{v-1}$ without proving that it annihilates
$\mathcal P_{v-1}$.  IQ5 uses only the expressly specified map on $W$.  A
conductor receiver from a larger source must be constructed on its original
observation image before using IQ5.

\section{Evaluation of the full inherited minimum}
For $m>0$ define
\[
 Q=(AG^{-1}A^*)^{-1},\qquad R_G=G^{-1}A^*Q,\qquad
 P_G=R_GA.                                          \tag{IQ6}
\]
The inverse exists: if $b\ne0$, surjectivity of $A$ makes $A^*b\ne0$,
so $b^*AG^{-1}A^*b>0$.  Direct multiplication gives
\[
 AR_G=I_Z,\quad P_G^2=P_G,\quad P_G^*G=GP_G,
 \quad\ker P_G=\ker A.                              \tag{IQ7}
\]
Every vector with $Az=b$ has a unique expression $z=R_Gb+n$, $n\in\ker A$.
Its cross term vanishes because
\[
 (R_Gb)^*Gn=b^*QA n=0,\qquad
 (R_Gb)^*G(R_Gb)=b^*Qb.
\]
Consequently the complete original minimum is evaluated, not merely bounded:
\[
 \begin{split}
 \|h\|_{\mathcal J_W}^2
   &=b^*Qb+\alpha^*H\alpha+n^*Gn,\\
 \min_{\mathcal Rh=b}\|h\|_{\mathcal J_W}^2
   &=b^*Qb
    =\min_{\Lambda w=b}\|w\|_W^2,\\
 s_W(b)&=\iota_-O^{-1}G^{-1}A^*Qb.
 \end{split}                                       \tag{IQ8}
\]
The minimum is unique: both omitted positive terms vanish exactly when
$\alpha=0$ and $n=0$.  The eight-state minimizing vector has coefficients
$\beta_j/2,-\beta_j/2$ at $f_{j,+},f_{j,-}$, respectively; IQ1 has not been
replaced by a Euclidean normalization.  For $m=0$ the only observation value
is zero, and its unique minimum is $h=0$.

For clarity, the Gram matrix in the original $(\alpha,\beta)$ coordinates is
\[
 M=\begin{pmatrix}
 H+E^*GE&E^*GO\\ O^*GE&O^*GO
 \end{pmatrix},\qquad
 M^{-1}=\begin{pmatrix}
 H^{-1}&-H^{-1}E^*O^{-*}\\
 -O^{-1}EH^{-1}&O^{-1}(G^{-1}+EH^{-1}E^*)O^{-*}
 \end{pmatrix}.                                    \tag{IQ9}
\]
These identities follow by multiplying the triangular coordinate change
$S=\left(\begin{smallmatrix}I&0\\E&O\end{smallmatrix}\right)$, its inverse
$S^{-1}=\left(\begin{smallmatrix}I&0\\-O^{-1}E&O^{-1}\end{smallmatrix}\right)$,
and $M=S^*\indmindiag(H,G)S$.  In particular the full evaluation row
$\mathsf R=A(E\ O)$ satisfies
\[
 \mathsf R M^{-1}\mathsf R^*=AG^{-1}A^*.              \tag{IQ10}
\]
Thus the cancellation of the even-label block follows from the actual
attained minimum; it is not a deletion of that block before minimization.
The complete observation kernel is
\[
 \ker\mathcal R=
 \{\iota_+\alpha+\iota_-O^{-1}(n-E\alpha):
             \alpha\in\indminC^4,\ n\in\ker A\},          \tag{IQ11}
\]
of dimension $8-m$.  It retains the four-dimensional evaluation graph
kernel ($n=0$) and all additional unseen target directions.

\section{The exact source analogue and conductor receiver}
Let $L=T_A|_U:U\to W$ be the original invertible four-plane conductor,
and retain the original maps $LX=Y$, $L\Phi=\Psi$.  Put
\[
 H_U=X^*\langle\ ,\ \rangle_UX,\qquad
 G_U=\Phi^*\langle\ ,\ \rangle_U\Phi,
 \qquad\mathcal J_Uh=(X\alpha,\Phi z).               \tag{IQ12}
\]
The source observation used here is expressly $\Lambda L:U\to Z$.
Its coordinate matrix is $\Lambda L\Phi=A$; it is not an unidentified
quotient of an ambient observation.  Repeating the proof IQ6--8 gives
\[
 Q_U=(AG_U^{-1}A^*)^{-1},\quad R_U=G_U^{-1}A^*Q_U,
 \quad s_U(b)=\iota_-O^{-1}R_Ub,\quad
 \min_{\Lambda Lu=b}\|u\|_U^2=b^*Q_Ub.              \tag{IQ13}
\]
The conductor $L\oplus L$ intertwines the two total maps and induces the
identity map on the common observation value $b\in Z$.  This identity has
source metric $Q_U$ and target metric $Q$; no equality of those metrics is
assumed.  Its two original minimizing lifts differ by the explicit vector
\[
 d(b)=L\Phi R_Ub-\Psi R_Gb
     =\Psi(R_U-R_G)b\in\ker\Lambda,\qquad
 \|L\Phi R_Ub\|_W^2=b^*Qb+\|d(b)\|_W^2.            \tag{IQ14}
\]
Indeed $A(R_U-R_G)=0$, and the orthogonality calculation in IQ8 applies.
If $\sigma_{\min},\sigma_{\max}$ are the actual extremal singular values
of $L$ for these original norms, then minimization of
$\sigma_{\min}^2\|u\|_U^2\le\|Lu\|_W^2
\le\sigma_{\max}^2\|u\|_U^2$ over the same fibre gives
\[
 \sigma_{\min}^2 Q_U\preceq Q\preceq\sigma_{\max}^2 Q_U. \tag{IQ15}
\]
This proves the exact metric receiver and its inherited bounds without
asserting a new estimate for the original conductor singular values.

\section{Full monodromy action and every observation defect}
For a signed root permutation, retain the matrices
$\Pi e_j=e_{\pi(j)}$ and $Re_j=\sigma_j e_{\pi(j)}$.
The source calculation SE16 gives
\[
 B_g=ORO^{-1},\quad C_g=E\Pi-B_gE,\qquad
 g:(\alpha,z)\longmapsto(\Pi\alpha,B_gz+C_g\alpha).
                                                               \tag{IQ16}
\]
Substitution of $z=E\alpha+O\beta$ proves this formula directly from
$(\alpha,\beta)\mapsto(\Pi\alpha,R\beta)$.  In particular it retains
the complete mixing block, and has the exact group law.  For
$z=R_Gb+n$, $n\in\ker A$, its observed value is
\[
 \mathcal R(gh)=T_gb+AB_gn+AC_g\alpha,\qquad
 T_g=AB_gR_G.                                         \tag{IQ17}
\]
This is the entire defect from an operation on the quotient observation.
It proves the exact criterion: $g$ induces a linear map on that quotient
if and only if
\[
 AB_g(\ker A)=0\quad\hbox{and}\quad AC_g=0.          \tag{IQ18}
\]
Necessity follows by applying IQ17 to the two independently variable
parts of IQ11; sufficiency follows by the same formula.  In that case
the induced map is exactly $T_g$, independently of the chosen representative.

Every two-sign flip actually occurs in the monodromy calculated in
SM10--14.  Indeed the square of the displayed adjacent four-cycle flips
the two adjacent signs.  The three vectors $(1,1,0,0),(0,1,1,0),(0,0,1,1)$
span all even sign vectors over $\mathbb F_2$, proving the assertion for
every pair, without treating an arbitrary signed permutation as a proved
geometric loop.

For the flip $D_J$ of a two-element set $J$, IQ16 becomes
\[
 B_J=OD_JO^{-1},\qquad
 C_J=O(I-D_J)F,
 \qquad \indminim C_J=O\indminSpan\{e_j:j\in J\}.              \tag{IQ19}
\]
To prove the last equality, projection of $\indminim F$ onto the coordinates in
$J$ is onto: prescribe those two coordinates, set one of the remaining
coordinates to zero, and solve for the other using its nonzero $a_k$ in
IQ2.  Since $I-D_J$ multiplies exactly the coordinates in $J$ by two,
the assertion follows with the scalar two retained.  Therefore
\[
 \indminim(AC_J)=\indminSpan\{AOe_j:j\in J\},\quad
 \indminrank(AC_J)=\dim\indminSpan\{AOe_j:j\in J\},\quad
 \sum_{|J|=2}\indminim(AC_J)=Z.                            \tag{IQ20}
\]
The final equality holds since $O$ is invertible and $A$ is onto $Z$.
In particular, if the observation is nonzero, some $AOe_j$ is nonzero,
and every pair containing that index gives $AC_J\ne0$.  Thus no nonzero
linear observation of evaluation alone carries the full even-sign
monodromy action.  This statement is stronger than a failure of one
chosen four-dimensional presentation: it is proved for every nonzero
linear receiving observation of that evaluation.

Here is an explicit zero-class witness, including the branch factors.
Choose $j\in J$ with $AOe_j\ne0$ and an index $k\notin J$.  Put
\[
 x=e_j-(a_j/a_k)e_k\in\indminim F,\qquad
 F\alpha=x,\qquad
 h_0=\iota_+\alpha-\iota_-F\alpha.                 \tag{IQ21}
\]
Existence of $\alpha$ follows from the proved image identity IQ2.
Then $z(h_0)=0$, but
\[
 \mathcal Rh_0=0,\qquad
 \mathcal R(g_Jh_0)=2AOe_j\ne0.                    \tag{IQ22}
\]
There is no unidentified cancellation or asymptotic term in this witness.

\section{The defect measured in the inherited quotient metric}
Restrict the full defect to the evaluation graph kernel in IQ11, namely
$h_0(\alpha)=\iota_+\alpha-\iota_-F\alpha$.  Its original norm is
$\|h_0(\alpha)\|^2=\alpha^*H\alpha$.  Its image under an even flip
has observed quotient norm exactly
\[
 \|\mathcal R(g_Jh_0(\alpha))\|_Q^2
   =\alpha^*C_J^*A^*QA C_J\alpha.                  \tag{IQ23}
\]
Hence the complete squared operator norm of this mixing defect, from
the inherited graph-kernel norm to the inherited quotient norm, is
\[
 \|\mathcal Rg_J|_{\ker\mathcal E}\|^2
 =\lambda_{\max}
   \bigl(H^{-1/2}C_J^*A^*QA C_JH^{-1/2}\bigr).       \tag{IQ24}
\]
The formula follows from the Rayleigh quotient after the invertible
substitution $v=H^{1/2}\alpha$.  Positivity of $Q$ proves that IQ24
is strictly positive exactly when $AC_J\ne0$.  In particular, IQ21
gives the explicit lower bound
\[
 \|\mathcal Rg_J|_{\ker\mathcal E}\|^2
 \ge \frac{4\,(AOe_j)^*Q(AOe_j)}{\alpha^*H\alpha}>0. \tag{IQ25}
\]
If the smallest possible denominator is desired, let
$V=\indminim F$ carry the restriction of the original coordinate Euclidean
form solely for defining the adjoint in this explicit solution, and set
$B=(FH^{-1}F^*)|_V:V\to V$.  This operator is positive definite on
$V$: for $v\in V$, $v^*Bv=\|H^{-1/2}F^*v\|^2$, and
$F^*v=0$ with $v\in\indminim F$ forces $v=0$.  Then
\[
 \alpha_*=H^{-1}F^*B^{-1}x,\qquad
 F\alpha_*=x,\qquad
 \alpha_*^*H\alpha_*=x^*B^{-1}x.                   \tag{IQ26}
\]
For any other preimage $\alpha_*+t$, $Ft=0$ and
$\alpha_*^*Ht=x^*B^{-1}Ft=0$; this proves the claimed minimizing
denominator.  The resulting original-metric lower bound in IQ25 is
therefore $4\|AOe_j\|_Q^2/(x^*B^{-1}x)$.  The auxiliary coordinate
adjoint in IQ26 changes no norm in IQ23--25.

The canonical minimizing section in IQ8 also has an exact action defect.
Since its $\alpha$ coordinate is zero, the difference between its
monodromy image and the minimizing lift of the new observation is
\[
 g s_W(b)-s_W(T_gb)
 =\iota_-O^{-1}(I-P_G)B_gR_Gb\in\ker\mathcal R.     \tag{IQ27}
\]
The same orthogonality as in IQ8 evaluates its full cost:
\[
 \|g s_W(b)\|_{\mathcal J_W}^2
 =b^*T_g^*QT_gb+
  \bigl((I-P_G)B_gR_Gb\bigr)^*G
       \bigl((I-P_G)B_gR_Gb\bigr).                 \tag{IQ28}
\]
These formulas do not imply that the group preserves the original
positive forms.  They compute its actual failure to preserve the
minimizing section, with all original cross entries retained.
Finally, since $B_gB_h=B_{gh}$, the matrices of these compressed
operations have the exactly evaluated composition defect
\[
 T_gT_h-T_{gh}
   =-AB_g(I-P_G)B_hR_G.                            \tag{IQ29}
\]
This expression can vanish for particular observations; for example a
four-dimensional invertible observation has $P_G=I$.  Even in that
case the independent mixing defect IQ23 can be nonzero, and IQ22
still proves that these compressed operations do not define the
action of the entire eight-state space modulo evaluation.  Equations
IQ17, IQ23 and IQ27 specify all three exact receiving maps required
to compare these presentations.

\section{Provenance and limits of the derivation}
The input source is the programme proof
\texttt{work/rh\_counterfactual\_20260913/fable\_conductor\_20260919/}
\texttt{FABLE\_TO\_ORIGINAL\_CONDUCTOR.tex}, specifically SE1--18 and
SM1--19.  The finite-dimensional proofs IQ1--29 above are supplied in
full rather than attributed to a human source not used here.  They
prove the exact attained metric transfer and the complete observation
defect.  They do not assert that the ambient arithmetic observation
descends through an unverified source quotient, that the nonlinear
cover changes an original conductor kernel, or that the positive
defect yields an RH sign estimate.

\endgroup

% END COMPLETE PARITY AND CONDUCTOR RECEIVERS
